\documentclass[12pt]{article}
\usepackage[margin=1.0in]{geometry}
\usepackage{amssymb}	% for \square check-box
\usepackage{hyperref}
\usepackage{enumitem}
\usepackage{graphicx}


\pagestyle{empty}
\newcounter{lastNum}   % allows for interspersing text between questions
\begin{document}
	
\begin{center}
	\textbf{Data Science Techniques (MAT 339)} \\
	\textbf{Homework 8} - - \textit{Initial batch of problems, more to come} -- \textbf{20 points}\\
\end{center}
By the beginning of class on \underline{Wednesday, April 1st}, turn in a \textbf{hard-copy} of answers to questions asking for hard-copy responses and \textbf{upload Python files} to Canvas for questions asking you to program TUIs, following the naming convention specified in each problem. \\

\begin{enumerate}
		
	\item Build a \textbf{counter TUI} using \texttt{textual}. The TUI should have one button to increase the counter and another to reset the counter. The counter should be displayed in the TUI. Each button should be given a clear, understandable label. Hint: you can use an instance variable to keep track of the counter. Name this file \textbf{\texttt{counter\_buttons.py}}.
	
	\item The goals of this question are for you to increase your familiarity with \texttt{textual} and to gain more experience "reading the docs", meaning going through the technical documentation of a software library on your own. Answer this question using the \texttt{textual} widget reference page at  			\href{https://textual.textualize.io/widgets/}{https://textual.textualize.io/widgets/}. \textbf{You may not use AI to help you answer this question.}  
	\begin{enumerate}
		\item For each of the widgets in the list below, find on the widget's reference page the \textbf{parameters} that can be passed into the widget when it is created (look for entries in the "Table of contents" on the right side of the screen that begin with \texttt{param}).  Choose \textbf{2-3} parameters you think might be useful to use and put them in a small table that has these columns: \underline{parameter name}, \underline{description} (paraphrase in your own words if possible), and \underline{example} (create an example value that could make sense). Do this for these widgets: \textbf{Input, Header, Link, Select}, and \textbf{1 more of your choice}.
		\item \textbf{Messages} in \texttt{textual} are triggered by \textbf{events}. This question is similar to the previous one but instead of finding \textbf{parameters}, this time find \textbf{messages} relating to widgets. See the "Messages" entry in the "Table of contents" on the right side of the screen. Create a table again, this time with the columns: \underline{message name} and \underline{meaning} (i.e. what does it mean when this event is generated). Fill the table with \textbf{at least 1 }message for each of these widgets: \textbf{Input, Button, Select}, and \textbf{1 more of your choice}.
		\item When an event triggers a message, a widget parameter may be changed. Looking over the parameters you found in part (a) or others from the documentation, create a table with \textbf{at least three} such parameters (so there will be at least 3 rows in this table) having the columns: \underline{name of parameter that changes}, \underline{parameter's associated widget}, \underline{event causing the change}. 
	\end{enumerate}

	\item Build a \textbf{TUI} using \texttt{textual} in a file named \texttt{divisor.py} that reports whether one user-inputted number (the "dividend") can be divided with no remainder by another user-inputted number (the "divisor"). The TUI should meet these requirements:
	\begin{itemize}
		\item accept each number from an \textbf{Input} widget
		\item include a label near each input widget identifying which number goes where (e.g. "dividend" or "divisor"); hint: consider a \textbf{Label} widget
		\item the above input and label widgets should fill a vertical column on the left-hand side of the TUI (hint use \texttt{with Vertical} after importing \texttt{Vertical()} from \texttt{textual.containers})
		\item there should be a right-hand side column containing a message to the user saying whether the divisor divides the dividend or not, e.g. "7 does divide 21", "4 does not divide 9". To simplify any error-handling, consider using a \texttt{try-except} block that tries to divide the inputs. If the division operation can occur, report one of the messages above. If not, the \texttt{except} part of the block should set a message saying "Input valid numbers" or similar.
	\end{itemize}
	\vspace{.1in}
	
	An example of this TUI and its behavior is shown below. 		
	
	\vspace{.1in}
	\includegraphics[width=2.7in]{divisor-y.png}\qquad
	\includegraphics[width=2.7in]{divisor-n.png}
	

	\item Build a \textbf{TUI} using \texttt{textual} in a file named \texttt{text\_view.py} that allows the user to open, view and change the appearance of text files in a text area portion of the TUI. The TUI should meet these requirements:
	\begin{itemize}
		\item has two vertical containers - the left one should have buttons labeled "Bold" and "Blue" and a Select widget with names of three files; the right one should be a Static widget 
		\item when the user selects a name from the Select widget, the contents of the file by that name should appear in the Static widget on the right side (they should be loaded from the file, not hard-coded into this py file)
		\item when the user clicks the "Bold" button, the text in the Static widget should turn bold
		\item when the user clicks the "Blue" button, the text in the Static widget should turn blue		
	\end{itemize}
	
	The appearance should resemble the screen shots below.
	
	\includegraphics[width=2.7in]{haiku1.png}\qquad
	\includegraphics[width=2.7in]{haiku2.png}
	
	\sloppy
	Three sample files (haikus) are provided in this homework's folder but you can use your own files.  The samples came from: \url{https://thehaikufoundation.org/omeka/files/original/50a61193bc1caf7a3c86b9acc96d4b44.pdf}
	
	Hint: you can add the \texttt{rich}-style formatting we covered to the beginning of the Static widget's content to set bold and blue color.
	
	\item \textit{One final problem will be added soon.}
	

\end{enumerate}
\end{document}

